Federated reinforcement learning enables household batteries to share control knowledge while retaining local operating records, but compromised participants can corrupt the shared policy. We evaluate FedGrid, an aggregation method that combines solar-capacity grouping, robust within-group filtering, norm clipping and reference-based weighting. The testbed couples residual Soft Actor-Critic (SAC) controllers with export enforcement and unbalanced power-flow simulation for 36 households. Six poisoning attacks are evaluated using Ausgrid profiles, with additional retraining and transfer experiments using the Australian Capital Territory (ACT) NextGen data. Across five Ausgrid seeds, the final configuration, which filters before clipping and increases the clipping bound, lowers observed mean annual Normalized-attack damage from \$88 to \$11 per household relative to the original configuration. The paired change is $-\$76.8\pm\$67.5$ under a descriptive 95\% Student-$t$ interval; the exact two-sided signed-rank test gives $p=0.0625$ before adjustment. Reward poisoning remains a weakness, with approximately \$112 mean annual damage in ACT retraining. In the clean FedAvg comparison, solar curtailment is approximately 69\% lower than under the self-consumption rule, whereas bill savings are modest. An analysis of aggregation coefficients shows that filter admission does not guarantee positive weight, qualifying the interpretation of low recorded rejection. Cross-partition ACT imputation and post-hoc configuration selection further limit the evidence. The results support attack-specific evaluation that distinguishes economic damage, filter admission and effective contribution.
